% Generated draft: requires mathematical and source review.
\documentclass[11pt]{article}
\usepackage[margin=25mm]{geometry}
\usepackage{fontspec}
\setmainfont{Latin Modern Roman}
\usepackage{amsmath,amssymb,mathtools,unicode-math}
\setmathfont{Latin Modern Math}
\usepackage{enumitem,needspace,xurl,hyperref}
\setlength{\parindent}{0pt}
\setlength{\parskip}{5pt}
\setlength{\emergencystretch}{3em}
\newcommand{\N}{\mathbb{N}}
\newcommand{\Z}{\mathbb{Z}}
\newcommand{\Q}{\mathbb{Q}}
\newcommand{\R}{\mathbb{R}}
\newcommand{\C}{\mathbb{C}}
\newcommand{\F}{\mathbb{F}}
\newcommand{\A}{\mathbb{A}}
\newcommand{\PP}{\mathbb{P}}
\newcommand{\E}{\mathbb{E}}
\newcommand{\eps}{\varepsilon}
\newcommand{\dd}{\,\mathrm{d}}
\begin{document}
\section*{UnsolvedMath ID: 20000382}
\textbf{Source problem code:} catalogue entry
\subsection{Definitions and mathematical statement}
The AIM question asks which definitions of automorphism groups give the ``nicest formulas.'' Neither the desired formulas nor a precise object category is specified. The attached report says the Principles (A) and (B) mentioned in neighboring questions are absent from the archived sources and does not reconstruct them. It works instead with finite field extensions, covers, and curves, subject to essentially finite groupoid and effective-descent hypotheses.

For an essentially finite groupoid \(\mathcal G\), its usual mass is
\[
|\mathcal G|=\sum_{[x]\in\pi_0\mathcal G}\frac1{|\operatorname{Aut}(x)|}.
\]
For a finite group \(G\) acting on a finite set \(S\), the quotient-groupoid formula is \(|[S/G]|=|S|/|G|\). The report calls inverse weighting by the full arithmetic stabilizer the unique quotient-compatible convention. A precise uniqueness statement also needs normalization, additivity, and the exact class of allowed quotient presentations; those axioms are not supplied. Arithmetic stabilizers mean automorphisms over the field of definition preserving the entire moduli object, not automatically the full geometric automorphism group.

The report offers independent diagnostic identities. For a finite separable extension \(L/K\), with permutation action \(\rho\) of the absolute Galois group on \(\Omega=\operatorname{Hom}_K(L,K^{\mathrm{sep}})\), the field-extension identity is
\[
\operatorname{Aut}_K(L)\cong C_{\operatorname{Sym}(\Omega)}(\rho(\operatorname{Gal}(K^{\mathrm{sep}}/K))).
\]
For a finite geometric automorphism group \(\Gamma\) with Frobenius action and effective forms, its asserted twist identity is
\[
\sum_{[\text{Frobenius twists }x]}\frac1{|\operatorname{Aut}_{\F_q}(x)|}=1.
\]
The sum must include all effective arithmetic twists of the fixed geometric object; a missing descent datum or a different selection of forms changes the statement.

For a cover \(f:C\to D\), the report distinguishes fixed-target automorphisms \(\{\alpha:f\alpha=f\}\), moving-target automorphisms \(\{(\alpha,\beta):f\alpha=\beta f\}\), and abstract-source automorphisms \(\operatorname{Aut}(C)\). Markings and field structure must also be preserved if part of the object.

Finally its candidate rigidification comparison is
\[
\frac{dM_t}{N_t}\longrightarrow1
\quad\Longleftrightarrow\quad
\text{the density of the extra-automorphism locus tends to zero}.
\]
The summary does not define \(d,M_t,N_t\) or the parameter \(t\). In a finite-object interpretation one would need a fixed baseline stabilizer order \(d\), mass \(M_t\) using arithmetic stabilizers, a compatible unweighted count \(N_t\), and a specified extra-stabilizer locus. In particular unweighted arithmetic objects, geometric classes, and coarse \(\F_q\)-points cannot silently be identified: twists can create multiple arithmetic forms above one coarse point. The all-twists mass-one identity makes this distinction essential. Thus the displayed equivalence is retained as a source claim requiring its exact counting/rigidification conventions, not a universal formula for every coarse moduli count.
\subsection{Short English statement}
Choose automorphism groups to match the actual moduli objects and quotient formulas. The report favors full arithmetic stabilizers, gives centralizer and twist-mass diagnostics, and proposes an asymptotic comparison whose coarse-count and rigidification conventions remain undefined.
\subsection{Sources}
\begin{itemize}
\item AIM workshop ``Arithmetic statistics over finite fields and function fields,'' January 27--31, 2014, Question 4: \url{https://aimath.org/pastworkshops/arithstatffieldproblems.pdf}. No PDF was accessed here.
\item Local canonical registry \texttt{aim-algebraic-number-theory-notes.json}, \texttt{notes[33]}, tag \texttt{question}; permanent UnsolvedMath ID \(20000382\), source code \texttt{catalogue entry}. Entire original question asks which automorphism-group definitions give the nicest formulas. Original remarks and literature field are empty.
\item Supplied research attempt 1 proposing quotient compatibility and the centralizer, twist-mass, and rigidification diagnostics. It explicitly says neighboring Principles (A) and (B) cannot be recovered. No authoritative full bibliographic references are supplied. The report is machine generated, not peer reviewed, and requires independent verification before citation.
\end{itemize}
\end{document}
